% Canonical node 002 · I002; current section path 1.1.2.
% Canonical owner: I002
\cnnaNodeHeading{002}{I002}{Finite approximant depth $L$}
\cnnaNodePlacement{1.1.2}{D0}

\paragraph*{Scientific role.}
\texttt{I002} is the second free structural input.  It fixes the terminal
provenance depth of the finite approximant.  It is neither a spatial
coordinate nor an event time, and it does not encode an infinite-limit
sentinel.

\paragraph*{Formal contract.}
The admissible domain is
\begin{equation}
  \mathcal L := \mathbb N_0,
  \qquad L\in\mathcal L.
  \label{eq:i002-depth-domain}
\end{equation}
Thus every nonnegative integer is admissible, including $L=0$.  The value is
supplied explicitly and no default is derived.  At this node alone, $L=0$
means only ``zero permitted provenance extension depth''; the later
root-only interpretation additionally uses the root construction
\texttt{C002} and the address grammar \texttt{C003}.

\paragraph*{Executable and formal encoding.}
Python uses a frozen \texttt{FiniteApproximantDepth} carrier and rejects
negative or non-built-in-integer inputs.  Lean uses a structure with one
\texttt{Nat} field.  No proof field is needed because \texttt{Nat} already
represents exactly the required nonnegative values.  The theorem
\texttt{zeroAdmissible} records explicitly that zero is an ordinary member of
the domain rather than a special missing-value marker.

\paragraph*{Boundary countercheck.}
$L=0$ and arbitrary positive integers are admitted; negative integers,
floating-point values, strings, \texttt{None}, and Boolean values are rejected
at the Python boundary.  No test or theorem in this node establishes a
continuum or infinite-depth limit.  Such a limit, if later constructed, is a
separate downstream problem.

\paragraph*{Result and handoff.}
The node yields one validated finite cutoff.  Edge \texttt{E002} supplies it to
\texttt{C003}; edge \texttt{E081} later fixes the completion depth of
\texttt{C019}; and proof-support edge \texttt{E148} makes the same finite bound
available to \texttt{P004}.  None of these handoffs changes the meaning of
$L$.

\noindent\textbf{Primary code anchors.}
Python: \path{derivation/code/python/src/cnna/derivation/s01_primitive_response_coupled_finite_provenance/s01_primitive_inputs_grammar_and_event_order/s02_i002__finite_approximant_depth_l.py},
\texttt{FiniteApproximantDepth}, lines 20--46.  Lean:
\path{derivation/code/lean/core/src/CNNA/Derivation/S01_PrimitiveResponseCoupledFiniteProvenance/S01_PrimitiveInputsGrammarAndEventOrder/S02_I002_FiniteApproximantDepthL.lean},
\texttt{FiniteApproximantDepth}, lines 18--23; \texttt{zeroAdmissible}, lines
29--34.  Exact hashes and tests are given in the Supplement.
